%1. What does Plan Cuadrante identify?
%The paper states in various places that permanent geographic assignment and greater local knowledge are the key mechanisms underlying the RF effects of the policy. I think this interpretation is stronger than what the evidence can establish.
%
%Plan Cuadrante appears to have changed several features of frontline policing simultaneously. Assigning officers to smaller quadrants changes not only the officer-area relationships, but also the spatial organization of police presence. Depending on quadrant size, number of officers, and deployment practices, the reform may have effectively changed police density and patrol intensity, as well as response times. Each of these margins could plausibly affect crime, victimization, perceptions of safety, and citizen cooperation. Greater local knowledge is one compelling mechanism within this bundle, but it is not separately observed from these other operational changes.
%
%I therefore think the paper should explicitly acknowledge that the results identify the effect of a broader organizational and deployment reform rather than the effect of local knowledge alone. A more defensible interpretation would be that Plan Cuadrante, combining permanent officer-to-quadrant allocation with an increase in police density, improved public safety and citizens' perceptions of the police. Even holding the total number of officers fixed, reorganizing officers across smaller territorial units may change the intensity and concentration of local police presence.
%
%One potentially useful exercise would be to characterize police density across quadrants, for example in terms of officers per square kilometer (or per capita). The illustrative map in the paper suggests that quadrants vary substantially in geographic size (and the smaller quadrants are the ones with higher baseline crime rates). If the authors observe the number of officers assigned to each quadrant, it would therefore be useful to document the distribution of officers per quadrant, and examine whether treatment effects vary with this measure of effective local police density. Police density may itself be correlated with response times, patrol intensity, and other operational margins, so it would still not isolate a single mechanism. But it would help characterize an important dimension of what the reform changed.


